\itemhead{第3章 小结与复习 · 第2题（教材p.89）}{XJ-B1-C03-SU-P089-Y1026}
\label{XJ-B1-C03-SU-P089-Y1026}
{\footnotesize 湘教版教材p.89 / PDF第96页 / 第3章 小结与复习 / 第2题}
\field{原题}{数学讲究表示。函数有哪几种常用的表示法？试结合具体实例，分析、比较这几种表示法的特点。}
\field{参考答案}{解析法便于精确求值和代数推理；列表法明确离散数据及对应；图象法直观显示变化和最值。三者应互相校验，图象读数通常受作图精度限制。}
{\small \field{关联知识}{F04 函数三种表示法、图象与读数（教材帮p.85、86、93、94 / PDF94、95、102、103）；F03 函数求值、解析式与函数建模（教材帮p.80、87、88、89、90、131 / PDF89、96、97、98、99、140）}
\field{方法与易错点}{把关键点与定义域一起标到图上，逐区间读交点和上下关系。 易错：将示意图预测当事实；把两曲线交点收进严格不等式；折线漏端点。}
\field{答案依据}{答案册未提供本栏目的独立对应答案；依据原题、教材概念补写参考回应。开放性总结不设唯一数值答案。}}
